\chapter{MODELAGEM UML BASEADA EM COMPONENTES DE SOFTWARE (CBSE)}

A modelagem arquitetural do compilador MiniPar 2026.2 adota formalmente o paradigma de Engenharia de Software Baseada em Componentes (\textit{Component-Based Software Engineering} -- CBSE), utilizando a notação padrão da Unified Modeling Language (UML). 

O sistema é concebido como um conjunto integrado de componentes desacoplados, de alta coesão e baixo acoplamento, nos quais a comunicação e a transferência de artefatos entre as fases de compilação ocorrem estritamente através de contratos bem definidos de interfaces fornecidas (\textit{provided interfaces}) e interfaces requeridas (\textit{required interfaces}).

\section{Diagrama de Casos de Uso}

O Diagrama de Casos de Uso da Figura \ref{fig:uml_casos_uso} modela as interações dos atores (Usuário, Desenvolvedor e Aluno Iniciante) com as três modalidades de interface do compilador.

\begin{figure}[htbp]
\centering
\begin{tikzpicture}[
    actor/.style={circle, draw, minimum size=0.8cm, inner sep=0pt},
    usecase/.style={ellipse, draw, align=center, fill=blue!5, minimum width=3.8cm, minimum height=0.9cm, font=\small},
    systembox/.style={draw=black!70, thick, rounded corners, inner sep=0.5cm}
]
    % Atores
    \node (user) at (-5, 0) [actor, label=below:{Desenvolvedor / Aluno}] {};
    \draw (user.north) -- ++(0,-0.8) -- ++(-0.3,-0.6) (user.north) ++(0,-0.4) -- ++(0.3,-0.6);
    \draw (user.north) ++(-0.4,-0.2) -- ++(0.8,0);

    % Sistema
    \node[usecase] (uc1) at (0, 3) {Fornecer Código Fonte (.minipar)};
    \node[usecase] (uc2) at (0, 1.5) {Inspecionar Fases Didáticas\\(Tokens, AST D3, Semântica)};
    \node[usecase] (uc3) at (0, 0) {Compilar Programa com/sem\\Otimização (-O)};
    \node[usecase] (uc4) at (0, -1.5) {Gerar e Baixar Binário (.exe)};
    \node[usecase] (uc5) at (0, -3) {Executar Programa com Entrada};

    % Fronteira do sistema
    \begin{scope}[on background layer]
        \node[systembox, fit=(uc1) (uc5), label=above:{Sistema Compilador MiniPar 2026.2}] {};
    \end{scope}

    % Relacionamentos
    \draw[->, thick] (user) -- (uc1);
    \draw[->, thick] (user) -- (uc2);
    \draw[->, thick] (user) -- (uc3);
    \draw[->, thick] (user) -- (uc4);
    \draw[->, thick] (user) -- (uc5);
\end{tikzpicture}
\caption{Diagrama de Casos de Uso do Compilador MiniPar 2026.2 (Notação UML)}
\label{fig:uml_casos_uso}
\end{figure}

\section{Arquitetura de Componentes de Software}

A especificação estabelece que a arquitetura do compilador deve ser representada por meio de um Diagrama de Componentes na notação UML, destacando as conexões de interfaces fornecidas (representadas pelo símbolo de esfera/pirulito) e interfaces requeridas (representadas pelo semicírculo/soquete).

A Figura \ref{fig:uml_componentes} ilustra a arquitetura CBSE implementada no diretório \texttt{src/}:
\begin{itemize}
    \item O componente Analisador Léxico consome o texto-fonte e fornece a interface \texttt{ILexico} com a operação \texttt{analise\_lexica()}.
    \item O componente Analisador Sintático requer \texttt{ILexico} e fornece a interface \texttt{ISintatico} com a operação \texttt{analise\_sintatica()}.
    \item O componente Tabela de Símbolos fornece a interface \texttt{ITabelaSimbolos} para gerenciamento de escopos e consulta.
    \item O componente Analisador Semântico requer \texttt{ISintatico} e \texttt{ITabelaSimbolos}, fornecendo a interface \texttt{ISemantico}.
    \item O componente Otimizador no backend requer \texttt{ISemantico} e fornece a interface \texttt{IOtimizador} com transformações na AST.
    \item O componente Gerador de TAC fornece \texttt{ITAC} para síntese do código intermediário.
    \item O componente Gerador de Código C++ fornece \texttt{ICodegen}, transpilando para C++20 e orquestrando o \texttt{g++}.
\end{itemize}

\begin{figure}[htbp]
\centering
\begin{tikzpicture}[
    component/.style={rectangle, draw=black, thick, fill=blue!10, minimum width=3.4cm, minimum height=1.3cm, align=center, font=\footnotesize\bfseries},
    interface/.style={circle, draw=black, thick, fill=white, minimum size=0.4cm, inner sep=0pt},
    socket/.style={arc, draw=black, thick},
    conn/.style={-{Latex[length=2mm]}, thick}
]
    % Nós dos Componentes
    \node[component] (lex) at (0, 6) {<<component>>\\C\_Lexico};
    \node[component] (syn) at (0, 4) {<<component>>\\C\_Sintatico};
    \node[component] (sem) at (0, 2) {<<component>>\\C\_Semantico};
    \node[component] (sym) at (5.5, 2) {<<component>>\\C\_TabelaSimbolos};
    \node[component] (opt) at (0, 0) {<<component>>\\C\_Otimizador};
    \node[component] (tac) at (0, -2) {<<component>>\\C\_TAC};
    \node[component] (gen) at (0, -4) {<<component>>\\C\_Codegen};

    % Interfaces e conexões (Pirulito e Soquete)
    % ILexico
    \draw[thick] (lex) -- (0, 5.15) node[interface, label=right:{\footnotesize\texttt{ILexico}}] (i_lex) {};
    \draw[conn] (syn) -- (i_lex);

    % ISintatico
    \draw[thick] (syn) -- (0, 3.15) node[interface, label=right:{\footnotesize\texttt{ISintatico}}] (i_syn) {};
    \draw[conn] (sem) -- (i_syn);

    % ITabelaSimbolos
    \draw[thick] (sym) -- (3.7, 2) node[interface, label=above:{\footnotesize\texttt{ITabelaSimbolos}}] (i_sym) {};
    \draw[conn] (sem) -- (i_sym);

    % ISemantico
    \draw[thick] (sem) -- (0, 1.15) node[interface, label=right:{\footnotesize\texttt{ISemantico}}] (i_sem) {};
    \draw[conn] (opt) -- (i_sem);

    % IOtimizador
    \draw[thick] (opt) -- (0, -0.85) node[interface, label=right:{\footnotesize\texttt{IOtimizador}}] (i_opt) {};
    \draw[conn] (tac) -- (i_opt);

    % ITAC
    \draw[thick] (tac) -- (0, -2.85) node[interface, label=right:{\footnotesize\texttt{ITAC}}] (i_tac) {};
    \draw[conn] (gen) -- (i_tac);

    % ICodegen
    \draw[thick] (gen) -- (0, -4.85) node[interface, label=right:{\footnotesize\texttt{ICodegen}}] (i_gen) {};
\end{tikzpicture}
\caption{Diagrama de Arquitetura de Componentes de Software (UML CBSE com Interfaces Fornecidas e Requeridas)}
\label{fig:uml_componentes}
\end{figure}

\section{Modelagem via Diagrama de Classes Conceitual}

O projeto contempla a modelagem dos componentes e estruturas de dados centrais em um Diagrama de Classes UML. A Figura \ref{fig:uml_classes} detalha a hierarquia dos componentes concretos e as classes que realizam os nós fundamentais da AST.

\begin{figure}[htbp]
\centering
\begin{tikzpicture}[
    conn/.style={-{Triangle[open]}, thick}
]
    % Interfaces
    \node (i_lex_c) at (-3.5, 3) [draw, thick, fill=gray!8, inner sep=0.1cm] {
        \begin{tabular}{l}
        <<interface>> ILexico\\
        \hline
        + analise\_lexica(codigo: str): list[Token]
        \end{tabular}
    };

    \node (c_lex_c) at (-3.5, 1) [draw, thick, fill=gray!8, inner sep=0.1cm] {
        \begin{tabular}{l}
        C\_Lexico\\
        \hline
        + analise\_lexica(codigo: str): list[Token]
        \end{tabular}
    };
    \draw[conn] (c_lex_c) -- (i_lex_c);

    \node (i_syn_c) at (2.5, 3) [draw, thick, fill=gray!8, inner sep=0.1cm] {
        \begin{tabular}{l}
        <<interface>> ISintatico\\
        \hline
        + analise\_sintatica(tokens: list[Token]): Program
        \end{tabular}
    };

    \node (c_syn_c) at (2.5, 1) [draw, thick, fill=gray!8, inner sep=0.1cm] {
        \begin{tabular}{l}
        C\_Sintatico\\
        \hline
        + analise\_sintatica(tokens: list[Token]): Program
        \end{tabular}
    };
    \draw[conn] (c_syn_c) -- (i_syn_c);

    % AST Nodes
    \node (ast_base) at (0, -1) [draw, thick, fill=gray!8, inner sep=0.1cm] {
        \begin{tabular}{l}
        <<abstract>> ASTNode\\
        \hline
        + line: int\\
        + col: Optional[int]\\
        \hline
        + to\_dict(): dict
        \end{tabular}
    };

    \node (fork_node) at (-4, -3.2) [draw, thick, fill=gray!8, inner sep=0.1cm] {
        \begin{tabular}{l}
        ForkStmt\\
        \hline
        + destino: Destino\\
        + body: ASTNode
        \end{tabular}
    };
    \draw[conn] (fork_node) -- (ast_base);

    \node (join_node) at (-1.5, -3.2) [draw, thick, fill=gray!8, inner sep=0.1cm] {
        \begin{tabular}{l}
        JoinStmt\\
        \hline
        + args: List[ASTNode]
        \end{tabular}
    };
    \draw[conn] (join_node) -- (ast_base);

    \node (class_node) at (1.5, -3.2) [draw, thick, fill=gray!8, inner sep=0.1cm] {
        \begin{tabular}{l}
        ClassDecl\\
        \hline
        + name: str\\
        + superclass: str\\
        + attributes: List[VarDecl]\\
        + methods: List[FuncDecl]
        \end{tabular}
    };
    \draw[conn] (class_node) -- (ast_base);

    \node (prog_node) at (4.5, -3.2) [draw, thick, fill=gray!8, inner sep=0.1cm] {
        \begin{tabular}{l}
        Program\\
        \hline
        + statements: List[ASTNode]
        \end{tabular}
    };
    \draw[conn] (prog_node) -- (ast_base);

\end{tikzpicture}
\caption{Diagrama de Classes Conceitual dos Componentes e da AST (Notação UML)}
\label{fig:uml_classes}
\end{figure}

\section{Diagrama de Atividades / Fluxo de Compilação}

O fluxo dinâmico de execução do compilador, ilustrado na Figura \ref{fig:uml_fluxo}, detalha o percurso do programa desde o código-fonte até a emissão e execução do executável nativo.

\begin{figure}[htbp]
\centering
\begin{tikzpicture}[
    startstop/.style={circle, fill=black, minimum size=0.4cm},
    activity/.style={rectangle, draw=black, thick, fill=blue!8, rounded corners, align=center, font=\small, minimum width=4cm, minimum height=0.8cm},
    decision/.style={diamond, draw=black, thick, fill=yellow!15, align=center, font=\scriptsize, inner sep=0.1cm},
    arrow/.style={-{Latex[length=2.5mm]}, thick}
]
    \node[startstop] (start) at (0, 7) {};
    \node[activity] (a1) at (0, 6) {Código-Fonte MiniPar (\texttt{.minipar})};
    \node[activity] (a2) at (0, 4.8) {Análise Léxica (\texttt{C\_Lexico})};
    \node[activity] (a3) at (0, 3.6) {Análise Sintática (\texttt{C\_Sintatico})};
    \node[activity] (a4) at (0, 2.4) {Análise Semântica e Escopos (\texttt{C\_Semantico})};
    \node[decision] (d1) at (0, 1.0) {Otimização\\Habilitada?};
    \node[activity] (a5) at (3.5, 1.0) {Otimização AST\\(\texttt{C\_Otimizador})};
    \node[activity] (a6) at (0, -0.4) {Geração de TAC (\texttt{C\_TAC})};
    \node[activity] (a7) at (0, -1.6) {Transpilação C++20 (\texttt{C\_Codegen})};
    \node[activity] (a8) at (0, -2.8) {Compilação Nativa via \texttt{g++ -O3}};
    \node[activity] (a9) at (0, -4.0) {Executável Nativo Final (\texttt{.exe})};

    \draw[arrow] (start) -- (a1);
    \draw[arrow] (a1) -- (a2);
    \draw[arrow] (a2) -- (a3);
    \draw[arrow] (a3) -- (a4);
    \draw[arrow] (a4) -- (d1);
    \draw[arrow] (d1) -- node[above]{\scriptsize Sim} (a5);
    \draw[arrow] (a5) |- (a6);
    \draw[arrow] (d1) -- node[right]{\scriptsize Não} (a6);
    \draw[arrow] (a6) -- (a7);
    \draw[arrow] (a7) -- (a8);
    \draw[arrow] (a8) -- (a9);
\end{tikzpicture}
\caption{Diagrama de Atividades do Fluxo de Compilação do MiniPar 2026.2}
\label{fig:uml_fluxo}
\end{figure}
